\documentclass[11pt]{article}

\usepackage[margin=1.15in]{geometry}
\usepackage{amsmath}
\usepackage{amssymb}
\usepackage{amsthm}
\usepackage{booktabs}
\usepackage{url}
\usepackage[hidelinks]{hyperref}

\newtheorem{theorem}{Theorem}[section]
\newtheorem{proposition}[theorem]{Proposition}
\newtheorem{lemma}[theorem]{Lemma}
\newtheorem{corollary}[theorem]{Corollary}
\theoremstyle{definition}
\newtheorem{definition}[theorem]{Definition}
\newtheorem{remark}[theorem]{Remark}

\newcommand{\art}[1]{{\texttt{\hyphenchar\font=`\-\relax #1}}}
\newcommand{\LB}{\mathrm{LB}}

\title{A Farkas certificate for a failed discharging argument}

\author{%
  (author list to be completed)%
}

\date{\today}

\begin{document}
\sloppy
\maketitle

\begin{abstract}
A discharging argument is parameterised by the amounts of charge its rules move.
Fix the rule shapes and the pool of reducible configurations; then the set of
amount vectors under which the argument closes is a \emph{polyhedron}, because
each local structure that the pool fails to block contributes one linear
inequality in the amounts. Failure of the argument is thus infeasibility of a
linear program, and infeasibility carries a Farkas certificate: a short
nonnegative combination of the inequalities that sums to a contradiction,
independent of any solver and checkable in exact rational arithmetic.

We give the certificate format, a solver-free exact-rational checker, and the
first instance. Fix the $84$ discharging rule shapes of the $2026$
near-linear-time proof of the Four Colour Theorem \cite{nl4ct2026} and restrict
its published $8{,}200$-configuration D-reducible pool to the $5{,}895$ members
of ring size at most $14$. A $12$-row irreducible certificate then shows that no
nonnegative amount vector $x\in\mathbb{R}^{84}_{\ge0}$ satisfies the local charge
conditions that schema imposes over that pool: the argument cannot be made to
close there by re-tuning its amounts. We say ``certify'' rather than ``prove''
in a sense made precise in \S\ref{sec:notclaim} --- the certificate is exact and
every row is re-derived from its source by the checker, but the row set is
necessary only under hypotheses that we state inside the theorem. The
contradiction is legible: it lives entirely at the degree-$5$/degree-$7$
interface, and $65$ of the $84$ rule shapes carry a zero coefficient in every row
of it. A positive control --- the same construction over the full pool --- is
feasible, with the published amounts satisfying every generated row.

The result bears on a question Steinberger \cite{steinberger2010} posed in $2009$
and explicitly declined to answer, whether an unavoidable set of D-reducible
configurations of ring size at most $14$ exists; we settle it for this pool and
these rule shapes, and no further. Automated-discharging frameworks
\cite{stolee2014,bousquet2024square,lavalicov2022} solve the same kind of linear
program inside a loop and treat infeasibility as a signal to enlarge the
configuration set; the certificate their loops discard is exactly the object we
report.

Two secondary results are included. We determine $f(r)$, the minimum number of
interior vertices of a D-reducible configuration of ring size $r$, for
$r=8,\dots,13$: $f = 5,5,6,7,7,8$, exhaustively over $13.2$ million generated
near-triangulations and tight at every ring. And we report an independent
re-verification of the published 4CT configuration corpus, which located $430$
configurations of the $2026$ pool that no stage of its pipeline uses and a
counting defect in the $1995$ \texttt{reduce.c} reference oracle.
\end{abstract}
%======================================================================
\section{Introduction}
\label{sec:intro}
%======================================================================

A discharging argument has two separable parts. There are the \emph{rule shapes}
--- the local patterns along which charge may move --- and there are the
\emph{amounts} those rules move. Almost all of the literature's attention, and
all of the difficulty of designing such an argument by hand, is in the shapes.
But once the shapes and the pool of reducible configurations are fixed, the
amounts are just numbers, and the condition ``every vertex ends with non-positive
charge'' is, at each local structure the pool fails to block, a single linear
inequality in those numbers. The set of amount vectors that make the argument
close is therefore a polyhedron.

This observation is not new; it is the engine of the automated-discharging
frameworks of Stolee \cite{stolee2014}, Bousquet, Deschamps, de Meyer and Pierron
\cite{bousquet2024square} and La and Valicov \cite{lavalicov2022}. What is new is
what we do with it when the polyhedron is \emph{empty}. Farkas' lemma
\cite{farkas1902} says that an infeasible linear system admits a certificate: a
nonnegative combination of its inequalities that sums to a contradiction. Applied
here, that certificate is a short, human-readable object which says
\emph{precisely which local structures, in what proportions, make the schema
impossible} --- and it is checkable in exact rational arithmetic by a program
that needs no solver, no floating point, and no trust in the enumeration that
generated the other rows.

\subsection{The result}

The Four Colour Theorem has been proved four times by machine
\cite{appel1977every,appel1977every2,rsst1997,steinberger2010,nl4ct2026} and once
in a proof assistant \cite{gonthier2008}. In every case the proof pairs an
unavoidable set of reducible configurations with a discharging argument
certifying unavoidability, and in every case that argument arrives with its
amounts already chosen and no account of why those amounts and not others.

We give the certificate format, an exact-rational checker for it, and the first
instance. The instance concerns the $2026$ near-linear-time proof of the Four
Colour Theorem \cite{nl4ct2026}, whose discharging schema has $84$ rule shapes
and whose pool has $8{,}200$ D-reducible configurations of ring size up to $18$.
Restrict that pool to its $5{,}895$ members of ring size at most $14$ and call the
result $P_{14}$. Theorem~\ref{thm:lp} then says: \emph{no} nonnegative amount
vector $x\in\mathbb{R}^{84}_{\ge0}$ satisfies the local charge conditions the
schema imposes over $P_{14}$. The argument cannot be rescued there by re-tuning
its amounts, and the reason is twelve rows long.

The certificate is legible in a way that an LP verdict is not. Its twelve rows
are nine cartwheels with a degree-$7$ hub and three degree-$5$ wheels; no
degree-$8$, $9$ or $10$ structure appears at all, even though the empirical
failure at the published amounts is spread across hub degrees $7$ to $10$. Every
one of the nine degree-$7$ hubs has a degree-$5$ spoke. Only $19$ of the $84$
rule shapes carry a nonzero coefficient anywhere in it, so $65$ of the shapes are
irrelevant to the obstruction. And the contradiction reads as arithmetic: the
degree-$5$ rows force certain amounts up, the degree-$7$ rows force the same
amounts down, and once $P_{14}$ stops blocking four particular cartwheels the two
demands are incompatible, by exactly $10\cdot460 - 10\cdot462 = -20$.

This bears on a question posed in $2009$. Steinberger \cite{steinberger2010} gave
the first unavoidable set of D-reducible configurations, at the cost of ring size
up to $16$, and wrote that he ``was not able to find an unavoidable set using
only configurations of ring-size $14$ or less, but [had] no opinion as to whether
such a set exists or not''. We settle a sharply scoped version: for the $2026$
rule shapes and the $2026$ pool, ring $\le 14$ is impossible, and the twelve rows
say why. We do not settle Steinberger's question, and \S\ref{sec:scope} states
exactly how far short we fall.

\subsection{Why the result type is new}
\label{sec:novelty}

Discharging has been automated in essentially this way three times. Stolee's
\textsc{adage} \cite{stolee2014} builds a linear program whose variables are
discharging amounts and whose constraints are configurations, and reports the
optimum as a density bound. Bousquet et al.\ \cite{bousquet2024square} build the
same kind of program for square colouring of planar graphs and wrap it in a loop:
solve, and if the optimum is insufficient, enlarge the configuration set and
solve again. La and Valicov \cite{lavalicov2022} automate the verification side
of the same procedure. Cranston and West's survey \cite{cranston2017discharging}
gives the hand version of the method these frameworks mechanise.

In all three frameworks, an insufficient LP is a signal to change the input.
Bousquet et al.\ enlarge $\mathcal{C}$; Stolee tries a different named rule set.
The dual object that says \emph{why} the present input cannot succeed is
computed by every simplex solve --- it is the Farkas ray the solver returns
alongside its infeasibility verdict --- and it is reported by none of them.
That object is a theorem, not a diagnostic. It is finite, it is small, it names
its own hypotheses, and it can be checked independently of the search that found
it. Reporting it converts a failed search into a statement about the schema.

\paragraph{The analogy with unsatisfiability proofs.} The practice we are
advocating is standard one field over. A SAT solver reporting \textsc{unsat} is
now expected to emit a machine-checkable proof, and large computer-assisted
results are accepted on the strength of those proofs rather than of the solver:
the determination of the packing chromatic number of the infinite square grid
\cite{subercaseaux2023packing} rests on a $34$-terabyte compressed DRAT proof,
and the empty hexagon number \cite{heule2024hexagon} on a reduction verified in
Lean. A Farkas certificate is the linear-programming analogue of an
unsatisfiability proof, and in this setting it is very much smaller --- twelve
rows, checkable by hand in principle. The discharging literature has not adopted
the practice.

The difference in setting matters too: those frameworks run their loops on
schemas being designed, where enlarging the configuration set is the natural
move on failure, whereas we run ours on a \emph{published, complete} schema whose
configuration set is the object under study and cannot be enlarged without
changing the claim.

\subsection{What ``certified'' means here, and what we do not claim}
\label{sec:notclaim}

We use two grades and label results with them. \textbf{Proven} means a
mathematical argument, in some cases with a finite exhaustive case analysis whose
completeness argument is given; Theorem~\ref{thm:fr} is of this kind.
\textbf{Certified} means a finite computation whose output is an independently
checkable certificate --- here, a Farkas certificate verified in exact rationals
by a program that re-derives every row from its source rather than trusting the
certificate's own numbers. Theorem~\ref{thm:lp} is certified, not proven, and we
say ``certify'' rather than ``prove'' throughout for it.

The distinction is not decorative, because three hypotheses are load-bearing and
are stated inside the theorem or in \S\ref{sec:scope}. (i) The amounts are
assumed nonnegative; this is used twice, and it is exactly load-bearing rather
than decorative, since the monotonicity argument fails on measured examples once
negative amounts are permitted. (ii) Rows are generated for every local structure
the pool fails to \emph{block}, which is what the released code does but is not
the same as ``every local structure that is realizable''; an unblocked but
geometrically impossible structure would contribute a spurious row, which is
anti-conservative for an infeasibility verdict. (iii) Three of the twelve rows
come from our reconstruction of the degree-$5$/$6$ half of the discharge, which
\cite{nl4ct2026} discharges in a single sentence by citation. \S\ref{sec:rows}
gives that reconstruction, its line-by-line audit against the pinned source, and
an exhaustive refinement check; but it remains a reconstruction, and without it
the whole system is feasible.

Finally, the result is scoped to \emph{these} $84$ rule shapes and \emph{this}
pool. Other D-reducible configurations of ring $\le14$ exist outside the
published pool, and adding or replacing rule shapes is not covered.

\subsection{Secondary results, and what is reported elsewhere}

Two further results are included because they are short and because the second
is what licenses the first. \S\ref{sec:fr} determines $f(r)$, the minimum number
of interior vertices of a D-reducible configuration of ring size $r$, for
$r=8,\dots,13$; this is a clean exhaustive theorem, and the first quantitative
refinement we know of to Birkhoff's $1913$ ring-size bound \cite{birkhoff1913}.
\S\ref{sec:verify} reports an independent re-verification of the published 4CT
corpus by a reducibility checker written from the RSST mathematics rather than
ported from its code, including a $5{,}692{,}937$-wheel differential against the
released C++, $430$ configurations of the $2026$ pool that no stage of its
pipeline uses, and a counting defect in the $1995$ \texttt{reduce.c} reference
oracle. The broader motivation is the one Tao \cite{tao2025machine} articulates
for machine-assisted proof: a machine proof becomes mathematics when it yields
transferable structure. A machine-learning line ran alongside this work and is
reported separately.
%======================================================================
\section{Preliminaries}
\label{sec:prelim}
%======================================================================

We keep this section to what the certificate needs and refer to
\cite{rsst1997,rsst-reduce} for everything standard.

\paragraph{Configurations and reducibility.} A \emph{configuration} $K$ of ring
size $r$ with $k$ interior vertices is given by its \emph{free completion}
$S(K)$: a simple plane near-triangulation of the closed disk on $n = r+k$
vertices whose outer face is bounded by an induced cycle $R$ of length $r$ (the
\emph{ring}) with every ring vertex of degree $\ge3$, whose interior is nonempty
and connected with every interior vertex of degree $\ge5$, and all of whose
bounded faces are triangles. RSST's reference parser \texttt{ReadConf} in
\texttt{reduce.c} records three further conditions on any file it accepts, one of
them redundant; the substantive one is that each interior vertex meets the ring
in at most two separate contiguous arcs. Let $\Phi(K)$ be the set of
$3$-colourings of $E(R)$, up to the $S_3$ action on colours, that extend to a
tri-colouring of $S(K)$, and write $a(K) = |\Phi(K)|$. Then $K$ is
\emph{D-reducible} if iterating the Kempe-closure operator from $\Phi(K)$ reaches
$\emptyset$, and \emph{C-reducible} if the same holds after contracting a
\emph{contract} $X$ of edges. The published catalogs record a nonempty $X$
exactly for those configurations that are not already D-reducible; reducibility
in the C-or-D sense is therefore strictly weaker than D-reducibility. The $2026$
pool \cite{nl4ct2026} and Steinberger's set \cite{steinberger2010} are D-only;
the RSST $633$ are not.

\paragraph{Discharging schemas.} A discharging argument assigns an
\emph{initial charge} $T_0(v)$ to each vertex of a hypothetical minimum
counterexample, moves charge along edges by a finite set of local \emph{rules},
and derives a contradiction from the fact that total charge is positive while
every vertex ends non-positive. RSST \cite{rsst-discharge} organise the local analysis around
\emph{axles}; Steinberger \cite{steinberger2010} and \cite{nl4ct2026} use
\emph{cartwheels} --- a hub together with its first and second neighbourhoods,
with degrees carried as intervals rather than exact values, refined by branching
until each is determined. A configuration in the pool \emph{blocks} a cartwheel
if it appears inside it; blocked cartwheels are discarded, since a minimum
counterexample contains no reducible configuration.

The $2026$ schema, which is the one we analyse, sets
$T_0(v) := 10\,(6-d(v))$, so that $\sum_v T_0(v) = 120$ by Euler's formula
(\cite{nl4ct2026}, source line $602$ and Lemma \texttt{lem:120}). Each rule $R$
is a planar degree-interval pattern with a distinguished dart; it sends an
\emph{amount} $r(R)$ of charge from a source vertex $s(R)$ to a target vertex
$t(R)$ whenever the pattern matches. The final charge is
\begin{equation}
\label{eq:finalcharge}
  T(u) \;=\; T_0(u) \;+\; \sum_{v \sim u}\bigl(\phi(v,u) - \phi(u,v)\bigr),
\end{equation}
where $\phi(v,u)$ is the total amount sent from $v$ to $u$. The published schema
fixes $r(R) \in \{1,2\}$ for all $84$ rules.

\begin{remark}[$43$ figures, $84$ rules]
\label{rem:84}
\cite{nl4ct2026} draws $43$ rules; each is counted twice, once per reflection,
except two that are reflection-symmetric, giving $43\cdot 2 - 2 = 84$ (source
line $635$: ``Figure \texttt{fig:rules} shows 84 rules since precisely two, the
first and fourth last, are symmetric under the reflection''). This matches the
$84$ files in \art{third\_party/discharging-rules/R} one for one. The paper also
notes that its $12$th and $13$th rules are a single rule in Steinberger's set of
$42$. Our amount vector ranges over all $84$ coordinates independently, which is
\emph{weaker} than requiring reflection pairs to share an amount --- the safe
direction for an infeasibility verdict.
\end{remark}

\paragraph{Survivors are permitted.} This is the point on which a naive encoding
of ``the discharging argument closes'' goes wrong, so we state it early. The
charge argument need not kill every unblocked local structure: the $2026$ proof
leaves $5{,}439 + 6{,}790 + 3{,}285 + 626 + 8 = 16{,}148$ stage-one wheels alive
on purpose, and three gluing lemmas dispose of the survivors. The correct local
condition is $T \le 0$, not $T < 0$. \S\ref{sec:controls} reports the negative
control that makes the difference concrete: the encoding demanding strict
negativity is infeasible even over the \emph{complete} $8{,}200$-configuration
pool, i.e.\ it would ``disprove'' the published theorem.
%======================================================================
\section{The amount polyhedron and its failure certificate}
\label{sec:lp}
%======================================================================

\subsection{Amounts as variables}

Fix a discharging schema: an initial charge $T_0$, a finite list of rule shapes
$R_1,\dots,R_N$, and a pool $U$ of reducible configurations. The schema's
\emph{amounts} are the numbers $r(R_1),\dots,r(R_N)$; collect them into a vector
$x \in \mathbb{R}^N_{\ge0}$. Everything else --- which patterns exist, which
cartwheels the pool blocks, how the refinement branches --- is fixed. The
question this paper asks is what the argument can do as $x$ varies.

The basic observation is that the final charge is affine in $x$. Let $L$ be a
local structure in which every degree is determined, so that for each rule shape
$R_k$ it is decided which darts at the hub $R_k$ fires along. Writing
$\mathrm{IN}_k(L)$ for the number of darts along which $R_k$ sends charge
\emph{into} the hub and $\mathrm{OUT}_k(L)$ for the number along which it sends
charge out, \eqref{eq:finalcharge} becomes
\begin{equation}
\label{eq:leafrow}
  C(L,x) \;=\; T_0(d) \;+\; \sum_{k=1}^{N}\bigl(\mathrm{IN}_k(L) - \mathrm{OUT}_k(L)\bigr)\,x_k ,
\end{equation}
an affine function with integer coefficients: no epigraph variables, no binaries,
no case split. The requirement that the hub of $L$ ends non-positive is one
linear inequality in $x$.

Consequently the set of amount vectors under which the argument closes over $U$
is an intersection of half-spaces with the nonnegative orthant:
\[
  X(U) \;=\; \bigl\{\, x \in \mathbb{R}^N_{\ge0} \;:\; A x \le b \,\bigr\},
\]
one row per local structure that $U$ fails to block. \emph{The schema closes for
some choice of amounts if and only if $X(U) \ne \emptyset$.} When
$X(U) = \emptyset$, Farkas' lemma \cite{farkas1902} supplies a certificate: a
vector $y \ge 0$ with
\begin{equation}
\label{eq:farkas}
  (y^\top\! A)_k \;\ge\; 0 \quad\text{for every } k, \qquad y^\top b \;<\; 0 ,
\end{equation}
whence for any candidate $x \ge 0$ the chain
$0 \le (y^\top\! A)\,x = y^\top(Ax) \le y^\top b < 0$ is a contradiction. The
certificate is a list of (row, multiplier) pairs. It is finite, it is checkable
by hand arithmetic in principle and by an exact-rational program in practice, and
--- this is the point of the format --- it names the small set of local
structures actually responsible. Nothing about it depends on the solver that
found it, or on floating-point arithmetic, or on trusting the enumeration that
produced the other rows.

\subsection{Where the rows come from}
\label{sec:rows}

Two families of rows are needed, and the second is not optional.

\paragraph{High-degree hubs.} The released $2026$ implementation enumerates hubs
of degree $7$ to $11$ and refines each cartwheel until every relevant degree is
determined. On a fully refined leaf $L$ all $d$ spokes are fixed, the per-spoke
maximisation inside \texttt{amount\_of\_possible\_charge\_send} never runs, and
the spoke combined rules are determined by the leaf graph, so \eqref{eq:leafrow}
applies verbatim. The invariant the code enforces at such a leaf yields the
necessary condition
\[
  C(L,x) \;\le\; 0, \qquad\text{and strictly } C(L,x) < 0 \text{ when either }
  d(L) \in \{9,10,11\}
\]
\[
  \text{or } d(L)\in\{7,8\} \text{ with no spoke of degree} \ge 7 ,
\]
the strict cases being those in which the leaf must be pruned before an assertion
is reached rather than surviving to the gluing lemmas. Note that $C = 0$ is
\emph{not} imposed: a leaf with $C<0$ is pruned earlier, so $C \le 0$ is exactly
the necessary-and-sufficient charge condition, weaker than what the code is
observed to satisfy at the published amounts. As a differential check on
\eqref{eq:leafrow}, all $10{,}094$ published full-pool bad cartwheels evaluate to
exactly $0$ under our independent implementation of the row semantics, and all
$10{,}094$ satisfy $d \in \{7,8\}$ with a spoke of degree $\ge 7$: the row
semantics reproduce all three C++ assertions on the published output with $0$
discrepancies.

\paragraph{Low-degree hubs.} Taken alone, the leaf rows are satisfied by
$x = 0$: send nothing, and every hub of degree $\ge 7$ keeps its negative initial
charge $10(6-d)$. Any encoding that stops there is vacuous, and we measured
exactly that (Table~\ref{tab:variants}). The missing constraint is at
$d \in \{5,6\}$, where $T_0 > 0$ and charge must actually leave. The $2026$ paper
does not verify that case: it asserts the stronger $T(v) = 0$ and discharges the
obligation in one sentence by citation (source line $1144$), and the cited
computation \cite{steinberger2010} was carried out at \emph{Steinberger's}
amounts. Once $x$ is re-tuned that citation no longer applies as proved, so we
impose the necessary condition directly. For each hub degree $d\in\{5,6\}$ and
each spoke-degree necklace $w$ that the pool does not block we add the row
\begin{equation}
\label{eq:lowdeg}
  \LB(w,x) \;=\; 10(6-d) \;-\; \sum_i \mathrm{POSS}_i(x) \;+\; \sum_j \mathrm{ALW}_j(x)
  \;\le\; T(v) \;\le\; 0 ,
\end{equation}
where $\mathrm{POSS}_i$ over-counts sends (it sums $x_k$ over base rules not
\texttt{never\_apply} out of spoke $i$) and $\mathrm{ALW}_j$ under-counts receipts
(base rules that \texttt{always\_apply} into the hub from spoke $j$). Both
directions are conservative, so $\LB$ is a genuine lower bound on the hub's final
charge, valid for any $x \ge 0$ at any level of refinement.

Two things make these rows trustworthy rather than a reconstruction we merely
hope is faithful. First, a quote-by-quote audit of \eqref{eq:lowdeg} against the
pinned \LaTeX{} source of \cite{nl4ct2026}: the initial charge, the total $120$,
the definition \eqref{eq:finalcharge}, and the requirement at $d\in\{5,6\}$
(which the paper states as $T(v)=0$, stronger than our $T(v)\le0$) each match a
located line, and rule $R(1)$ decoded from the paper's own figure source is a
single edge with source of degree exactly $5$ and amount $2$, matching
\art{rule001.rule} byte for byte --- so the row $10 - 5x_{001} \le 0$ is forced.
Second, a sanity anchor: at the published amounts $x_0$ the maximum of $\LB$ over
all $3{,}046$ unblocked low-degree wheels is exactly $0$ and never positive. Had
we got this side wrong in the loose direction, $x_0$ would have violated it.

\paragraph{The coarse-refinement gap.} The low-degree rows are built on the
coarse wheel (hub plus spoke-degree necklace, second neighbourhood left open),
and blocking is tested there. ``Refining only strengthens the bound'' does not
cover the danger, because refining can also \emph{delete} the row: if a coarse
wheel is unblocked but every refinement of it is blocked, no vertex of a minimum
counterexample carries that necklace and the row is spurious --- which is
anti-conservative for an infeasibility verdict. So it was checked twice. Exact
model counting over the full refinement space of the certificate's degree-$5$
necklaces leaves between $35.0\%$ and $61.2\%$ of all refinements unblocked in
every case: not a knife edge. Independently, the gap closes uniformly for all
$86$ distinct low-degree rows at once, because
\texttt{PseudoConfiguration::representative\_degree} collapses any degree upper
bound above $8$ to the single value $9$, so the coarse blocking test \emph{is}
the blocking test of the tail-maximised refinement; all $580$ unblocked $d=5$ and
all $2{,}466$ unblocked $d=6$ necklaces have an unblocked tail-maximised
refinement, with $0$ exceptions of $3{,}046$.

\subsection{Why rows observed at one amount vector constrain every amount vector}

The refinement forest is explored with pruning, so which leaves are observed
depends on the amounts used to explore. The rows are nevertheless valid for all
$x \ge 0$.

\begin{lemma}[Monotonicity]
\label{lem:mono}
For every leaf $L$ observed during the $x_0$-pruned search and every $x \ge 0$,
the stated condition $C(L,x) \le 0$ (resp.\ $< 0$) is necessary.
\end{lemma}

\begin{proof}
(a) The branching operations \texttt{update\_degree\_by\_rule},
\texttt{concrete\_degree\_except\_tail}, \texttt{refine\_always} and
\texttt{refine\_never} only narrow degree intervals and never change rotations,
so the shape of the forest does not depend on $x$. It does depend on the pool,
because \texttt{fix\_in\_rules} branches over the non-blocked combined-rule set;
but that dependence moves in the safe direction, since a smaller pool blocks
fewer combined rules. We verified the containment concretely: the full-pool
combined rules are a content-exact subset of the $P_{14}$ ones
($671 \subset 681$ files), so the full-pool forest is a sub-forest of the
$P_{14}$ forest.
(b) \texttt{blocked\_by\_reducible\_configuration} is independent of $x$ and
monotone in the pool, and \texttt{prune\_by\_non\_associated\_rule} is
independent of both.
(c) \texttt{upper\_bound\_of\_charge} is non-increasing along refinement whenever
$x \ge 0$: narrowing a degree interval can only add \texttt{always\_apply}
matches, which raises the out-charge, and can only remove candidates surviving
\texttt{never\_apply}. This is the step that uses nonnegativity, and it is the
step whose faithfulness to the source is argued rather than derived; it was
additionally measured on $2{,}145$ parent/child pairs, where it holds with
$x \ge 0$ and fails on $8$ pairs once negative amounts are allowed.
(d) Hence for any $x\ge0$, either every ancestor of $L$ survives, in which case
the assertion at $L$ demands $C(L,x) \le 0$; or some ancestor $A$ has
$C(A,x) < 0$, and then $C(L,x) \le C(A,x) < 0$ by (c).
\end{proof}

Since $P_{14} \subseteq P_{\mathrm{full}}$ (verified: $0$ of the $5{,}895$ names
absent from the $8{,}200$), every full-pool leaf row is also valid for the
ring-$\le14$ problem, which is what lets a certificate mix rows of both origins.

\begin{remark}[Direction of the relaxation]
\label{rem:direction}
The row set omits the three gluing lemmas and omits every leaf the pruned search
never reached. Both omissions make the encoded feasible set \emph{larger} than
the true one. An \textsc{infeasible} verdict therefore transfers to the true
problem, while \textsc{feasible} is only a non-refutation. This asymmetry is why
the certificate, and not the LP verdict, is the reportable object.
\end{remark}

\subsection{Code-necessary versus math-necessary}
\label{sec:necessity}

One objection nearly killed the result and is generic to this method. The
released code counts only \texttt{always\_apply} rules outward, and stops
refining once no rule is \texttt{dominantly\_apply}-but-not-\texttt{always\_apply};
since \texttt{dominantly\_apply} is strictly stronger than
\texttt{has\_intersection}, a leaf can retain rules that are neither
\texttt{always\_apply} nor \texttt{never\_apply} on an outgoing dart --- rules
that might fire in some realization but are not counted in the out-charge.
Measured on the certificate's own leaves, $11$ of $5{,}292$ (rule, out-dart)
pairs are undetermined in this sense, and for those leaves $C$ strictly
\emph{over}-estimates the hub's true final charge, so $C(L,x)\le0$ would be a
condition of the released \emph{verification procedure} rather than of the
mathematics. Two blunt repairs --- dropping such rows, or weakening them ---
both restore feasibility (Table~\ref{tab:variants}), so the distinction is not
hypothetical.

The repair we use instead is to pass to the \emph{tail-maximised} leaf, narrowing
every open tail $[a,9]$ with $a<9$ to $[9,9]$; since $9$ abbreviates ``$9$ or
more'', this restricts attention to the sub-family of realizations in which every
unpinned second neighbour has large degree. If the tail-maximised leaf (i) has a
fully determined outgoing side, (ii) yields the identical row --- same
coefficients, same constant, same strictness --- and (iii) is still unblocked by
the pool, then in those realizations the out-charge is exact while the in-charge
can only exceed the branch's flag sum, so the true final charge is at least
$C(L,x)$ and $T\le0$ forces $C(L,x)\le0$. Undetermined incoming rules are
harmless precisely because they push the true charge up. All twelve rows of the
certificate below pass this audit: the nine degree-$7$ rows satisfy (i)--(iii)
individually, including the five whose original leaves had undetermined outgoing
pairs, and the three degree-$5$ rows are mathematically necessary by
construction, being already in lower-bound form
(\art{results/steinberger/iis\_math\_necessity.json}).
%======================================================================
\section{The instance: the ring-$\le14$ sub-pool of the $2026$ schema}
\label{sec:thm}
%======================================================================

\subsection{The question, and an empirical precursor}

Steinberger's unavoidable set \cite{steinberger2010} is D-only but reaches ring
size $16$; RSST's $633$ stay at ring $\le14$ but use contracts. Steinberger
raises the obvious question and declines to take a position
(\art{third\_party/arxiv-0905.0043/src/4c.tex}, line $137$):
\begin{quote}\small
``Appel and Haken used only configurations of ring-size at most $14$\dots\ The
unavoidable set of Robertson et al.\ also has only configurations of ring-size
$14$ or less. Our unavoidable set uses configurations of ring-size up to
$16$\dots\ \textbf{We were not able to find an unavoidable set using only
configurations of ring-size $14$ or less, but we have no opinion as to whether
such a set exists or not.}''
\end{quote}
The $2026$ pool \cite{nl4ct2026} is D-only, has $8{,}200$ configurations of ring
size up to $18$, and contains $5{,}895$ of ring size $\le 14$; call that sub-pool
$P_{14}$.

Before any linear programming we ran the released pipeline with the pool replaced
by $P_{14}$, the $84$ rule shapes and the published amounts $x_0$ left untouched.
It fails, and in a structured way: $681$ non-blocked combined rules rather than
the baseline $671$ (ten revived, all ten predicted in advance by a gap analysis
with $0$ unexpected re-blockings), a per-wheel job count of $16{,}157$ rather
than $16{,}148$, and then $833$ of those jobs aborting --- $780$ on
\texttt{assert(C == 0)} ($530$ at $d=7$, $250$ at $d=8$) and $53$ on
\texttt{assert(d == 7 || d == 8)} ($51$ at $d=9$, $2$ at $d=10$), the two failure
kinds perfectly stratified by hub degree.

That is a statement about $x_0$ alone, and the obvious response is to re-tune the
amounts. The theorem below says no re-tuning can work. The $833$ aborted jobs are
also the raw material for it: re-run against an \texttt{NDEBUG} build with the
assertions compiled out, they yield $77{,}185$ additional fully refined leaves
that the assert-guarded run never recorded.

\subsection{The theorem}

\begin{theorem}[Certified; see \S\ref{sec:notclaim} for the grade and \S\ref{sec:scope} for the hypotheses]
\label{thm:lp}
Fix the $84$ discharging rule shapes of \cite{nl4ct2026}
(\art{third\_party/discharging-rules/R}), and let $P_{14}$ be the sub-pool of the
published $8{,}200$-configuration D-reducible pool consisting of its $5{,}895$
members of ring size at most $14$. Then there is no vector of nonnegative rule
amounts $x \in \mathbb{R}^{84}_{\ge0}$ satisfying the local charge conditions
that the $2026$ schema imposes over $P_{14}$: that is, no $x \ge 0$ under which
every fully refined cartwheel leaf $L$ that $P_{14}$ fails to block satisfies
$C(L,x) \le 0$ --- strictly $C(L,x)<0$ in the cases listed in \S\ref{sec:rows}
--- and every unblocked degree-$5$ or degree-$6$ wheel $w$ satisfies
$\LB(w,x)\le 0$. In particular the discharging argument of \cite{nl4ct2026}
cannot be made to close over $P_{14}$ by re-tuning its amounts.
\end{theorem}

All twelve rows of the certificate are non-strict, so infeasibility holds over
the nonnegative reals, hence a fortiori over $\mathbb{Z}^{84}_{\ge0}$ and over
the box $\{1,2\}^{84}$ that \cite{nl4ct2026} allows itself.

The system behind the theorem has $48{,}509$ distinct rows: $97{,}860$ fully
refined cartwheels yield $48{,}423$ distinct leaf rows after deduplication by
(coefficients, constant, strictness), to which the $86$ distinct low-degree rows
are added. The leaves are $10{,}094$ published full-pool bad cartwheels,
$10{,}581$ from the $15{,}324$ per-wheel jobs that completed, and the $77{,}185$
recovered from the $833$ jobs that aborted. By hub degree the rows are
$40{,}257$ at $d=7$, $8{,}070$ at $d=8$, $93$ at $d=9$, $3$ at $d=10$, $26$ at
$d=5$ and $60$ at $d=6$; the $93$ and $3$ are exactly the strict rows
corresponding to the $51+2$ degree-range failures above, and the $3{,}046$
unblocked low-degree wheels collapse to only $86$ distinct rows.

\begin{table}[t]
\centering\small
\caption{The four headline variants and the two lossy diagnostics of
\S\ref{sec:necessity}. ``$x_0$ viol.'' counts rows the published amounts violate.
Rows are deduplicated by (coefficients, constant, strictness). Full table of
eight variants, and reproduction instructions, in the artifact repository
(see the availability statement).}
\label{tab:variants}
\begin{tabular}{llrrl}
\toprule
pool & low-degree rows & rows & $x_0$ viol. & status \\
\midrule
full $8{,}200$ (control) & included & 3{,}511 & \textbf{0} & \textbf{feasible} \\
full $8{,}200$ (control) & omitted  & 3{,}425 & 0 & feasible (vacuously, $x=0$) \\
$P_{14}$ (main)          & included & 48{,}509 & 9{,}647 & \textbf{infeasible} \\
$P_{14}$ (main)          & omitted  & 48{,}423 & 9{,}647 & feasible (vacuously, $x=0$) \\
\midrule
$P_{14}$, rows dropped   & included & 42{,}121 & 8{,}699 & feasible --- too lossy \\
$P_{14}$, rows weakened  & included & 52{,}400 & 8{,}741 & feasible --- too lossy \\
\bottomrule
\end{tabular}
\end{table}

\subsection{The certificate}
\label{sec:cert}

A solver returns a $19$-row Farkas support with $y^\top b = -30$. Greedy deletion
filtering shrinks it to the $12$-row \emph{irreducible infeasible subsystem} of
Table~\ref{tab:cert}: deleting any single one of the twelve restores
feasibility, verified for all twelve. This is the smallest human-inspectable
object carrying the whole result. Its arithmetic fits on one line: the nine
degree-$7$ rows have right-hand side $10$ and multipliers summing to $460$, the
three degree-$5$ rows have right-hand side $-10$ and multipliers summing to
$462$, so
\[
  y^\top b \;=\; 10\cdot 460 \;-\; 10\cdot 462 \;=\; -20 \;<\; 0 ,
\]
while $(y^\top\! A)_k \ge 0$ holds in all $84$ columns --- in fact $78$ of the
$84$ columns cancel exactly and only six carry slack, which is why the
contradiction is tight.

\begin{table}[t]
\centering\small
\caption{The $12$-row irreducible certificate for Theorem~\ref{thm:lp}, with
multipliers $y$ and $y^\top b = -20$. ``rhs'' is the row's right-hand side and
$C(x_0)$ the leaf charge at the published amounts, so a positive entry marks a
row $x_0$ violates. Every row is non-strict. The full $84$-column coefficient
vectors, the $19$-row support, and the machine-readable certificates are in the
artifact repository (see the availability statement).}
\label{tab:cert}
\begin{tabular}{rrlrrll}
\toprule
$y_i$ & $d$ & spoke degrees & rhs & $C(x_0)$ & source & origin \\
\midrule
75  & 7 & 5,5,8,5,6,5,7 & 10 & 0 & \texttt{d7\_1155\_1} & full $\cap$ $P_{14}$ \\
75  & 7 & 5,5,5,8,5,5,7 & 10 & 0 & \texttt{d7\_166\_1}  & full $\cap$ $P_{14}$ \\
25  & 7 & 5,5,7,5,5,7,7 & 10 & 0 & \texttt{d7\_710\_10} & full $\cap$ $P_{14}$ \\
75  & 7 & 5,5,7,5,6,5,8 & 10 & 0 & \texttt{d7\_722\_3}  & full $\cap$ $P_{14}$ \\
55  & 7 & 5,5,7,5,7,5,7 & 10 & 0 & \texttt{d7\_740\_19} & full $\cap$ $P_{14}$ \\
60  & 7 & 5,7,5,7,6,7,7 & 10 & \textbf{2} & \texttt{d7\_3358\_2133} & \textbf{$P_{14}$ only} \\
15  & 7 & 5,7,5,7,6,7,7 & 10 & \textbf{2} & \texttt{d7\_3358\_2134} & \textbf{$P_{14}$ only} \\
55  & 7 & 5,7,5,7,7,6,7 & 10 & \textbf{2} & \texttt{d7\_3372\_1881} & \textbf{$P_{14}$ only} \\
25  & 7 & 5,7,5,7,7,7,7 & 10 & \textbf{3} & \texttt{d7\_3376\_1132} & \textbf{$P_{14}$ only} \\
75  & 5 & 5,5,7,6,7     & $-10$ & $-2$ & \texttt{d5-55767} & low-degree \\
168 & 5 & 5,7,5,7,7     & $-10$ & 0    & \texttt{d5-57577} & low-degree \\
219 & 5 & 5,7,6,7,7     & $-10$ & 0    & \texttt{d5-57677} & low-degree \\
\bottomrule
\end{tabular}
\end{table}

\paragraph{What it says.} The certificate is legible in a way LP output usually
is not.
\begin{itemize}\itemsep2pt
  \item \textbf{No degree-$8$, $9$ or $10$ row appears.} The contradiction needs
        none of the $53$ degree-range failures nor any of the $250$ degree-$8$
        failures seen empirically. The irreducible cause sits entirely at hub
        degree $7$.
  \item \textbf{Every one of the nine degree-$7$ hubs has a degree-$5$ spoke},
        so the whole obstruction lives at the degree-$5$/degree-$7$ interface.
  \item \textbf{Only $19$ of the $84$ rule shapes carry a nonzero coefficient
        anywhere in the twelve rows} --- \texttt{rule001} together with the
        reflection pairs of \texttt{rule002}, \texttt{003}, \texttt{004},
        \texttt{008}, \texttt{009}, \texttt{010}, \texttt{011}, \texttt{023} and
        \texttt{032}. The remaining $65$ shapes are irrelevant to the
        obstruction: re-tuning them cannot help.
  \item \textbf{Read as arithmetic}, the degree-$5$ rows say a degree-$5$ vertex
        must give away all $10$ units of its charge, forcing the amounts on
        \texttt{rule001}--\texttt{rule004} \emph{up}; the degree-$7$ rows cap what
        a degree-$7$ hub with degree-$5$ neighbours may receive back, pulling the
        same amounts \emph{down}. With $P_{14}$ no longer blocking four offending
        degree-$7$ cartwheels, the two demands cannot both be met.
\end{itemize}

\paragraph{The infeasibility is caused by the ring-$\le14$ rows.} Delete the four
$P_{14}$-only rows from the $19$-row support and the remaining $15$ --- each of
which is also valid for the full pool --- are feasible, with $x_0$ satisfying all
$15$. The contradiction is not lurking in the full-pool constraints; it is
created by exactly the cartwheels $P_{14}$ fails to block. Independently, those
four rows come from wheels \texttt{d7\_3358}, \texttt{d7\_3372} and
\texttt{d7\_3376} (\texttt{d7\_3358} contributes two leaves) --- precisely three
of the wheels whose job the released C++ aborted with
\texttt{Assertion failed: (C == 0)}. Our row semantics and the C++ assertions
agree wheel by wheel on the certificate's own support.

\subsection{Checking it}

\art{tools/verify\_farkas.py} is solver-free and uses exact rational arithmetic
throughout; there is no floating point anywhere on the checking path.
Critically, it does not trust the certificate's stored coefficients: it
re-derives every row from its source --- cartwheel files re-parsed and re-run
through the leaf-row constructor, low-degree tags rebuilt from the necklace and
re-run through the lower-bound constructor --- and fails loudly on any mismatch.
It then checks $y \ge 0$, $y \ne 0$, $(y^\top\!A)_k \ge 0$ for all $84$ columns
and $y^\top b < 0$, prints the contradiction chain of \eqref{eq:farkas}, and
independently confirms which rows $x_0$ violates. Deliberate corruption of a
single coefficient is rejected with an exact claimed-versus-re-derived diff. A
reader who accepts the row semantics of \S\ref{sec:rows} need trust nothing else
in this paper to check Theorem~\ref{thm:lp}.

\subsection{Controls}
\label{sec:controls}

\paragraph{Positive control.} The identical construction over the full
$8{,}200$-configuration pool is feasible, and the published amounts $x_0$ satisfy
every one of its $3{,}511$ rows with $0$ violations
(Table~\ref{tab:variants}). So the encoding is not vacuously infeasible: it
distinguishes the two pools, and it accepts the published proof.

\paragraph{Negative control.} The encoding a naive reading suggests ---
\emph{every} unblocked stage-one wheel must have strictly negative charge --- is
not a correct rendering of ``the argument closes'', because survivors are
permitted (\S\ref{sec:prelim}). We ran it anyway. It is infeasible at iteration
$0$ on the \emph{full} pool, with a verified $10$-row certificate
($y^\top b = -175$; three degree-$5$, two degree-$6$ and five degree-$7$ rows).
That encoding would therefore ``disprove'' the published theorem, which is the
concrete demonstration that it is the wrong one. We report it as a measurement
and never as evidence for the headline.

\paragraph{Ablations and robustness.} The degree-$6$ rows together with all leaf
rows are feasible; the degree-$5$ rows alone are already infeasible. Six solver
configurations all report infeasible, and a uniform slack of $0.082$ on every
right-hand side is the minimum needed to restore feasibility. The formulation was
also handed to a reviewer prompted to refute it across ten named attack lines,
which produced one real error (an over-strong pool-independence claim in
Lemma~\ref{lem:mono}(a)), one scope correction (\S\ref{sec:necessity}), and one
genuine gap, closed exactly (\S\ref{sec:rows}) --- none of which changed the
certificate.

\subsection{Scope}
\label{sec:scope}

\begin{enumerate}\itemsep2pt
  \item \textbf{These $84$ shapes only.} Nothing here rules out a different or
        larger set of rule shapes; that requires column generation over planar
        rule patterns (\S\ref{sec:open}).
  \item \textbf{This pool only.} $P_{14}$ is exactly the ring-$\le14$ part of the
        published pool. Many D-reducible configurations of ring $\le14$ lie
        outside it.
  \item \textbf{The $d\le6$ rows are a reconstruction} (\S\ref{sec:rows}) and are
        load-bearing: without them the system is feasible.
  \item \textbf{Unblocked $\ne$ realizable.} Rows are generated for every local
        structure the pool fails to block, exactly as the released code decides
        which cartwheels to analyse. A structure that is geometrically impossible
        for a reason outside the pool would contribute a spurious row, which is
        anti-conservative for an infeasibility verdict. It is also precisely the
        standard the published proof uses.
  \item \textbf{$x\ge0$ is a hypothesis}, used in Lemma~\ref{lem:mono}(c) and in
        \eqref{eq:lowdeg}, and it is load-bearing rather than decorative
        (monotonicity fails on $8$ measured pairs once negative amounts are
        allowed). In the other direction there is nothing to hedge: the published
        schema fixes $r(R)\in\{1,2\}$, and
        $\{1,2\}^{84} \subset \mathbb{Z}^{84}_{\ge0} \subset
        \mathbb{R}^{84}_{\ge0}$, so deciding infeasibility over the nonnegative
        reals is strictly stronger than the schema's own amount space requires.
\end{enumerate}
\section{The verification that licenses the certificate}
\label{sec:verify}

None of the linear-programming argument of the preceding sections means anything unless the
reducibility checker feeding it, and the row semantics built on top of that checker, are correct.
Our checker \texttt{fourcolor.reduce.check} was written directly from the mathematics of
\cite{rsst-reduce} (the \texttt{reduce.tex} write-up, in particular its expansion of
\cite[Theorem~(3.2)]{rsst1997}), not ported from \texttt{reduce.c}. Agreement with the released C
programs is therefore genuine differential evidence, not a tautology: two independent readings of
the same theorem, implemented by different people at different times in different languages, are
being asked whether they compute the same numbers. Table~\ref{tab:verify} summarizes the five
legs of that comparison; full logs, configuration files, and rerun scripts are in the artifact
repository (see the availability statement).

\begin{table}[t]
\centering
\small
\caption{The five verification legs. ``Outcome'' reports what was certified computationally, not
proved mathematically. Row~3 has no published $|\Phi(K)|$ to compare against -- the released 2026
\texttt{.conf} files carry only an ``$n\ r$'' header -- so ``confirmed'' there means only that our
checker independently certifies each configuration D-reducible, not that it matches a prior count.}
\label{tab:verify}
\begin{tabular}{@{}p{0.20\textwidth}p{0.32\textwidth}p{0.42\textwidth}@{}}
\toprule
What & Scale & Outcome \\
\midrule
Catalog re-verification, RSST 633 (\cite{rsst-reduce} \texttt{unavoidable.conf}) &
633 configurations, rings 6--14 &
Computed $|\Phi(K)|$ and $|\mathrm{consistent}|$ match published header fields for all 633; 249
D-reducible, 384 C-reducible via a nonempty published contract and not D-reducible. \\
Catalog re-verification, Steinberger $\mathcal{U}_{2822}$ &
2,822 configurations, rings 6--16, all header $b=0$ &
2,822/2,822 agree with published headers, 0 disagreements. \\
2026 pool &
7,697 configurations &
Independently confirmed D-reducible, 0 anomalies, 12.3 CPU-hours. Honest accounting: $8{,}200 =
7{,}697$~verified $+\ 430$~unused-and-removed $+\ 73$~unverified (rings 17--18, outside our
engines' ring budget). \\
Wheel-level differential, Python reimplementation vs.\ released C++ &
5,692,937 wheel candidates, hub degrees 7--11 ($11{,}165 + 48{,}915 + 217{,}045 + 976{,}887 +
4{,}438{,}925$) &
Every degree reports \texttt{exact\_match: true}: 0 false positives, 0 false negatives. \\
End-to-end reproduction of released C++ pipeline &
8,200-configuration pool (19,754 after symmetry expansion), 84 rule shapes, Debug build with all
assertions live, wall 4h38m &
Every count matches the published proof: 1,832 combined rules, 671 non-blocked; bad wheels
$5{,}439/6{,}790/3{,}285/626/8$ at $d{=}7,\dots,11$; 10,094 bad cartwheels (9,366 at $d{=}7$, 728 at
$d{=}8$); three gluing lemmas all passed (5,365 / 1,618 / 298 cartwheels). \\
\bottomrule
\end{tabular}
\end{table}

\paragraph{Why the differential matters.} The 5,692,937-wheel exact match is what licenses
treating the leaf charge as an affine function of the amount vector: it is the bridge between our
Python row semantics and the C++ that the published proof actually runs, so the LP rows we hand
the solver are demonstrably the rows the proof means.

\subsection{Two defects found in released code}

\paragraph{A legacy oracle that aborts.} \texttt{build/reduce\_stein}, compiled from the
Steinberger-extended \texttt{reduce.c} (needed because the unmodified 1995 code has
\texttt{MAXRING}~$=14$), processes 783 of Steinberger's 2,822 configurations and then aborts with
\texttt{SIGABRT} (exit 134), truncating mid-number. The crashing input is preserved as
\art{build/crash783.conf}: a ring-16 configuration with $n=28$ and $a=139{,}480$. Our checker
completes it in 227.7~s, returning $a=139{,}480$, matching the published header. The binary
corroborates the first 783 and then dies; our own checker, run against the published headers,
carries the remaining 2,039 unaided.

\paragraph{A counting defect in \texttt{reduce.c}.} On 70 sampled exhaustively-generated
RSST-legal configurations run through \texttt{build/reduce\_rsst}, there were 0 verdict
mismatches but 6 oracle-inconclusive cases: on those 6 the oracle recomputed a different
$|\Phi(K)|$ than our checker supplied and aborted with \texttt{*** ERROR: DISCREPANCY IN NUMBER OF
EXTENDING COLOURINGS ***} before printing a verdict. A third, from-scratch count of $|\Phi(K)|$,
sharing no code path with the production checker (different edge indexing, different triangle
enumeration, no gauge fixing, no greedy variable ordering), agrees with our checker in all 6 and
disagrees with the oracle in all 6: (ring, $n$, ours~=~brute force, oracle) equal $(8, 12, 72,
48)$, $(9, 13, 147, 94)$, $(10, 15, 367, 230)$, $(10, 15, 365, 243)$, $(10, 15, 408, 250)$, and
$(11, 17, 995, 532)$. The evidence points to \texttt{reduce.c}'s \texttt{strip()} edge-numbering
heuristic -- a hand-optimized greedy ordering routine, documented in its own comments as chosen to
shrink the search tree -- as the outlier on a class of inputs (small, exhaustively generated,
generic near-triangulations) that its original validation against a hand-curated catalog never
exercised. None of the six configurations' D-reducibility verdict is in question: all are
non-D-reducible, and the oracle never disputed the verdict, only the intermediate count. Our
checker agrees with the oracle's header $a$ on all 3,455 hand-curated catalog configurations
(633~+~2,822) with zero mismatches, so this is not a general defect. See
\art{results/theorem/fr\_table/crosscheck.json}.

\subsection{A pool 430 configurations larger than it needs to be}

Coverage is sharply two-level. At the wheel level only 85 of the 8,200 configurations ever act as
a blocker anywhere across $d=7,\dots,11$. At the cartwheel level 7,422 distinct configurations
appear as blockers over 2,253,185 logged blocking events -- 90.5\% of the pool, and a lower bound,
since the instrumentation records only the first matching configuration per event. The intuition
that a small core carries the proof is right at the wheel stage and wrong at the cartwheel stage.

A first prune tracked three blocking roles and found 673 unused; removing them made the pipeline
fail (all three gluing checks aborted, the bad-cartwheel count collapsed from 10,094 to 0). That
failure discovered a fourth blocking role: the gluing lemmas also consult the pool (2,199 distinct
glue-blockers). Tracking all four roles gives 430 configurations with zero usage anywhere.
Removing exactly those and re-running the entire pipeline gives PASS with counts identical to the
full-pool proof in every field.

\begin{proposition}[Certified]
The 2026 proof goes through verbatim with a pool of 7,770 configurations rather than 8,200, a
reduction of 5.2\%.
\end{proposition}

This is a floor we reached, not a claimed minimum.

We are not aware of any prior independent verification of the 2026 pool's reducibility; the
released \texttt{computer-checks} repository (\art{third\_party/GIT\_PINS.txt}, commit
\texttt{6cb85666}) consumes the pool as an input and contains no reducibility checker of its own.
\section{The minimum interior size of a D-reducible configuration}\label{sec:fr}

Write $f(r)$ for the minimum number of interior vertices of a D-reducible configuration of ring size $r$.

\begin{theorem}[Proven, exhaustive]\label{thm:fr}
For each ring size $r \in \{8,\dots,13\}$, every D-reducible configuration of ring size $r$ has at least $f(r)$ interior vertices, where
\[
\bigl(f(8),f(9),f(10),f(11),f(12),f(13)\bigr) = (5,5,6,7,7,8).
\]
Each bound is tight: a D-reducible configuration with exactly $f(r)$ interior vertices exists in the published catalogs.
\end{theorem}

Birkhoff \cite{birkhoff1913} proved that a reducible configuration has ring size at least $6$; to our knowledge Theorem~\ref{thm:fr} is the first published lower bound on interior size as a function of ring size. Published catalogs give per-ring interior minima, but a catalog minimum is only an upper bound on $f(r)$ -- the exhaustive negative side, that nothing smaller can possibly be D-reducible, is what is new.

\subsection{The enumeration}

For each $r$ and each $n$ from $r+1$ to $r+f(r)-1$ we generated every near-triangulation of the disk with outer face of size $r$ on $n$ vertices using \texttt{plantri -P$r$ -c3 -a $n$} (\texttt{plantri} 5.5 \cite{brinkmann2007}, pinned with its SHA-256 digest), filtered to the RSST configuration class, and ran the D-reducibility checker on every survivor. \texttt{plantri} is a canonical-construction-path generator producing exactly one representative per isomorphism class, which is what makes ``sweep every $n$ below threshold'' an exhaustive search rather than a sample; with \texttt{-P} and no \texttt{-m}, \texttt{-c3} defaults to \texttt{-c3m3}, i.e.\ no chords and no outer-face vertex of degree $2$.

\begin{table}[h]
\centering
\begin{tabular}{@{}lccccccc@{}}
\toprule
$r$ & interior range & raw & cand & excl & valid & D-red. \\
\midrule
8 & 1--4 & 436 & 18 & 6 & 12 & 0 \\
9 & 1--4 & 666 & 35 & 17 & 18 & 0 \\
10 & 1--5 & 13{,}338 & 221 & 127 & 94 & 0 \\
11 & 1--6 & 293{,}555 & 1{,}631 & 1{,}028 & 603 & 0 \\
12 & 1--6 & 509{,}288 & 4{,}092 & 3{,}004 & 1{,}088 & 0 \\
13 & 1--7 & 12{,}389{,}976 & 36{,}392 & 27{,}888 & 8{,}504 & 0 \\
\midrule
total & & 13{,}207{,}259 & 42{,}389 & 32{,}070 & 10{,}319 & 0 \\
\bottomrule
\end{tabular}
\caption{Exhaustive sweep below the interior threshold for each ring size. ``raw'' is the count of near-triangulations from \texttt{plantri}; ``cand'' passes the interior filter (nonempty connected interior of minimum degree $5$); ``excl'' additionally fails the RSST ring-contact-arc condition; ``valid'' is the true RSST configuration class. Every D-reducibility verdict is $0$. \art{results/theorem/fr\_table/manifest\_r*.json}.}
\label{tab:frsweep}
\end{table}

RSST's reference parser \texttt{ReadConf} in \texttt{reduce.c} records three further conditions on any file it accepts -- one of them redundant -- and we enforce all three; the substantive one is that each interior vertex meets the ring in at most two separate contiguous arcs. This is genuine structure, not bookkeeping: it is not implied by the near-triangulation conditions, and its exclusion rate rises from $33\%$ at $r=8$ to $77\%$ at $r=13$. Excluded graphs were nevertheless retained with a D-reducibility verdict computed anyway, so a D-reducible-but-not-RSST-legal graph would have been surfaced loudly; none was -- every one of the $32{,}070$ exclusions is non-D-reducible too, so the filter never changed a verdict, only the population count. The $r=13$ cells at interior $6$ ($n=19$) and interior $7$ ($n=20$) were run $8$-way sharded by residue class, roughly $2{,}250$\,s and $13{,}300$\,s per shard, giving $1{,}445$ and $6{,}710$ valid configurations and $0$ D-reducible.

\subsection{Verification and tightness}

Up to $20$ RSST-legal below-threshold configurations per ring were serialised and run through \texttt{build/reduce\_rsst}, compiled from the pinned \texttt{reduce.c}: $70$ sampled, $0$ verdict mismatches, $6$ oracle-inconclusive. The $6$ inconclusive cases trace to a defect in the oracle's colouring count, adjudicated in our favour by a third from-scratch implementation; see \S\ref{sec:verify} for the numbers.

Tightness witnesses, verified both by our checker and by the oracle, come from the published catalogs: $r=8$: \texttt{D0003.conf} ($n=13$, $a=100$); $r=9$: \texttt{D0006.conf} ($n=14$, $a=211$); $r=10$: \texttt{D0015.conf} ($n=16$, $a=536$); $r=11$: \texttt{D0024.conf} ($n=18$, $a=1{,}363$); $r=12$: \texttt{D0143.conf} ($n=19$, $a=2{,}863$) -- all five from the 2026 D pool \cite{nl4ct2026} -- and $r=13$: \texttt{126.504366} from the RSST 633 \cite{rsst-reduce} ($n=21$, $a=6{,}954$). \art{results/theorem/fr\_table/tightness\_witnesses.json}.

Theorem~\ref{thm:fr} is scoped to D-reducibility. The corresponding minimum over configurations that are C-reducible via a nonempty contract but not D-reducible is strictly smaller at rings $8$ and $11$, so the D-scoping in the statement is not cosmetic.
%======================================================================
\section{Discussion and open problems}
\label{sec:open}
%======================================================================

\paragraph{Certificates for the automated-discharging frameworks.} The immediate
consequence of \S\ref{sec:lp} is not about the Four Colour Theorem. Any of the
frameworks of \cite{stolee2014,bousquet2024square,lavalicov2022} can emit a
Farkas certificate at no extra cost, because their solvers already compute one
whenever a round of the loop fails; what is missing is the decision to record it,
the row-provenance bookkeeping that makes it re-derivable, and an exact-rational
checker. The value is highest exactly where those loops currently give the least
information: when a configuration set is believed to be complete and the LP still
fails, the certificate distinguishes ``the rule shapes are wrong'' from ``this
particular local structure is not covered'', and it names the structure. We would
like to see one reported for a schema someone is actively designing.

\paragraph{Removing the fixed-shapes hypothesis.} Theorem~\ref{thm:lp} fixes the
$84$ published rule shapes; the natural strengthening --- ``no discharging schema
of bounded rule complexity closes over the ring-$\le14$ pool'' --- requires column
generation. Given a violated row, one must price out a new planar
degree-interval pattern with a charge dart that would repair it, and that pricing
subproblem is a search over planar patterns rather than a linear program. It is
the main technical obstacle. A tractable intermediate is to enumerate all rule
shapes below a fixed size bound and re-run the same machinery over the enlarged
column set: infeasibility there would be a substantially stronger statement, and
the certificate format is unchanged.

\paragraph{Steinberger's question proper.} Settling whether \emph{some} D-only
unavoidable set of ring size $\le14$ exists needs both of the above plus a way to
range over pools. We note only that the obstruction we found is not diffuse: it
is a single interface, and $65$ of the $84$ shapes play no part in it. If there
is a positive answer, the twelve rows say where new rules would have to act.

\paragraph{$f(r)$ beyond $r=13$.} The $r=13$ interior-$7$ cell alone required
$11{,}534{,}001$ generated graphs and roughly $30$ core-hours; $r=14$ at interior
$8$ is about two orders of magnitude larger and is out of reach of the present
pipeline on a workstation. The promising route is to prune the generation with a
necessary condition on $a(K)$ rather than to run the reducibility check on every
candidate. The corpus already gives the upper bounds $f(14)\le9$, $f(15)\le10$
and $f(16)\le10$ from catalog witnesses.

\paragraph{A coloring-mass law, in progress.} The pruning condition just
mentioned is the subject of a companion line of work, reported separately: an
exact closed form $a = (2^r + 2(-1)^r)/6$ for the wheel, a proven cap
$a \le 11\cdot2^{r+k-7}$ tight at every wheel, a proof that the growth constant
is $3/2$ at $k=2$, and an obstruction theorem showing that no bound factoring
through $a \le P(S,4)/24$ can reach the empirically sharp constant $4/3$. Paired
with an empirical threshold law it would reproduce the whole $f(r)$ table
analytically to within one interior vertex at every ring; the sharp cap is at
present a conjecture and the threshold half is corpus-supported only. The
current state, including the corpus and the frontier data, is in the artifact
repository (see the availability statement).

\paragraph{Formalisation.} Both results here are within reach of a proof
assistant given Gonthier's formalisation \cite{gonthier2008} of the reducibility
machinery. The finite content of Theorem~\ref{thm:fr} is a replay of $10{,}319$
configuration checks. Theorem~\ref{thm:lp} is more tractable still at the
arithmetic level --- twelve rows over the rationals --- but the row-derivation
step, from a cartwheel to its $84$ coefficients, carries the real formalisation
burden, and it is the step on which the theorem's faithfulness to
\cite{nl4ct2026} depends.

%======================================================================
\section*{Data and code availability}
\addcontentsline{toc}{section}{Data and code availability}
\label{sec:availability}
%======================================================================

The certificate (both the $19$-row support and the $12$-row irreducible core, in
machine-readable form with full $84$-column coefficient vectors), the
solver-free exact-rational checker \art{tools/verify\_farkas.py}, the row
generator \art{tools/lp\_discharge.py}, the enumeration manifests and
per-configuration records behind Theorem~\ref{thm:fr}, the catalog
re-verification reports, the pipeline reproduction logs and the pool-coverage
measurements are all in the artifact repository at
\url{https://github.com/<repository>/four_color_digestion}, which will be
archived with a DOI on publication. All third-party inputs --- the four arXiv
e-prints, the \texttt{plantri} generator, the released rule and configuration
files and the released C++ checker --- are pinned in-repository with recorded
SHA-256 digests (\art{third\_party/CHECKSUMS.sha256},
\art{third\_party/GIT\_PINS.txt}). Following Steinberger's precedent for the
$2{,}822$-configuration set, the bulk machine output --- the $48{,}509$ generated
rows, the $97{,}860$ refined cartwheels, the $13.2$ million generated
near-triangulations and their verdicts --- is not reproduced here in any form; it
is in the repository, with a claim-to-artifact table naming the file behind every
number printed above.

\bibliographystyle{plain}
\bibliography{refs}

\end{document}
